\scanpage{60}
Cela signifie que l'on a une kyrielle de suites $\Sigma$-exactes (donc exactes)
\[
\begin{aligned}
0\longrightarrow\Psi^0\longrightarrow \mathbb{C}^0\longrightarrow\Psi^1\longrightarrow0&\\
\hspace{2em}0\longrightarrow\Psi^1\longrightarrow \mathbb{C}^1\longrightarrow\Psi^2\longrightarrow0&\\
\hspace{4em}0\longrightarrow\Psi^2\longrightarrow \mathbb{C}^2\longrightarrow\Psi^3\longrightarrow0&\\
&\cdots
\end{aligned}
\]
Par hyp on sait que
\[
\Psi^i_{\bullet}=k\ \text{placé en dim }i
\]
% slanted
\margin{\unsure{c'est peut-être pas a)|n'est-ce pas a)}}
On trouve aisément les \deleted{\unsure{Hom}$(\Phi$} \inserted{hom} entre les $\deleted{\Phi}\mathbb{C}^i$, et \deleted{\unsure{Hom}$(\,$les} $\Psi^j$
\[
\deleted{\Phi}\mathbb{C}^i\longrightarrow\Psi^j,\quad\Psi^i\longrightarrow\deleted{\Phi}\mathbb{C}^j,\quad\Psi^i\longrightarrow\Psi^j\qquad(\deleted{\Phi}\mathbb{C}^i\longrightarrow\deleted{\Phi}\mathbb{C}^j\ \text{déjà connu})
\]
on peut résumer la situation en disant que sur la sous-catégorie pleine $P_1\subset P$ formée des objets isom. à des sommes de $\Phi^i$ et de $\Psi^j$, le foncteur $L\longmapsto L_{\bullet}=\operatorname{Hom}(\Phi^{\bullet},L)$ est pl. fidèle. Son image essentielle est formée des complexes de chaînes $L_{\bullet}$ finis tels que les composantes \inserted{$L_i$, les $Z_i$ et les $B_i$} \deleted{\unsure{$H_i(L_{\bullet})$}} soient libres de \unsure{r. fi.|t. f.}

(On peut voir que le \deleted{\unsure{couple}} \inserted{\unsure{syst.}} $(P_1,\mathbb{C}^{\bullet},\Psi^0,\varepsilon,\Sigma\unsure{\cap|d}P_1)$ est universel pour les cat. $k$-additives munies d'une \deleted{\unsure{un}} \inserted{classe} $\Sigma$ de suites courtes, et d'une $\Sigma$-résolution d'un objet \deleted{$\Psi^0$}$Z^0$ \unsure{--|\ldots}).

On va maintenant supposer de plus

d) Pour tte suite \inserted{$\Sigma$-}exacte de $P$
\[
0\longrightarrow L^{\unsure{\flat|\bullet|0}}\longrightarrow M\longrightarrow N\longrightarrow0
\]
% slanted
\margin{il suffit d'avoir l'un des deux, p. ex. le \unsure{premier|récipr.}}
et tout hom $L\longrightarrow L'$ (resp $N'\longrightarrow N$) la suite \deleted{\unsure{som}} déduite par somme amalgam. (resp. produit fibré) existe dans $P$
% archive sheet number: 59
